% Capitulo: Construir a própria aplicação de master data e governança
\chapter{Construir a própria aplicação de master data e governança}
\label{cap:aplicacao-propria}

\rascunho{
  \textbf{Ideia do capítulo.} Um guia prático, porque há pouco material sobre isso: o
  que considerar para uma empresa construir a própria aplicação de master data e
  governança, agora que ficou mais fácil. Não é um tutorial de uma tecnologia; é a lista
  de decisões que eu gostaria de ter tido antes de começar.
  \begin{itemize}
    \item Parte do que vem antes: master data puxado por um projeto de negócio
      (\veja{sec:master-data-puxado}) e o steward como revisor
      (\veja{sec:steward-tutor}).
    \item Elementos do capítulo: tabela de construir ou comprar
      (Tabela~\ref{tab:construir-comprar}), arquitetura mínima
      (Figura~\ref{fig:aplicacao-propria}), um trecho de código, e uma lista de
      verificação no fim (Tabela~\ref{tab:checklist-aplicacao}).
  \end{itemize}
}

% ---------------------------------------------------------------------------
\section{Construir ou comprar}
\label{sec:construir-comprar}

\rascunho{
  \begin{itemize}
    \item A IA baixou muito o custo de construir; não baixou o custo de manter. A pergunta
      que decide é quem vai cuidar da aplicação daqui a dois anos.
    \item Construir não precisa ser para sempre: começar com uma aplicação própria num
      domínio e migrar para uma ferramenta de mercado quando a escala pedir é um caminho
      legítimo, desde que os dados saiam fácil.
  \end{itemize}
}

\pendente{escrever}

\begin{table}[htbp]
  \centering
  \caption{Quando construir e quando comprar uma solução de master data e governança.}
  \label{tab:construir-comprar}
  \small
  \begin{tabularx}{\textwidth}{@{}>{\raggedright\arraybackslash}p{0.17\textwidth}
      >{\raggedright\arraybackslash}X>{\raggedright\arraybackslash}X@{}}
    \toprule
    Critério & Construir faz sentido quando & Comprar faz sentido quando \\
    \midrule
    Escopo & Um ou dois domínios, com regras muito próprias do negócio & Muitos domínios e fontes, com regras comuns de mercado \\
    \addlinespace
    Equipe & Há quem programe e quem mantenha depois & Não há quem mantenha código interno \\
    \addlinespace
    Regulação & Um histórico de mudanças simples basta & Auditoria formal, certificações, setor regulado \\
    \addlinespace
    Prazo & É preciso algo em semanas, um domínio por vez & É um programa corporativo de longo prazo \\
    \addlinespace
    Custo & A licença é desproporcional ao problema & Manter código próprio custaria mais que a licença \\
    \addlinespace
    Integração & É preciso encaixar em fluxos internos específicos & A ferramenta já conecta com toda a pilha de dados \\
    \bottomrule
  \end{tabularx}
  \fonte{elaborado pelo autor.}
\end{table}

% ---------------------------------------------------------------------------
\section{Começar pelo problema, não pela plataforma}
\label{sec:comecar-pequeno}

\rascunho{
  \begin{itemize}
    \item Um domínio, uma regra, um dono: \enquote{clientes com nomes diferentes no ERP e
      no CRM, revisados pelo time comercial}. A plataforma nasce do problema, não o
      contrário.
    \item O primeiro usuário é o projeto de negócio que puxou o domínio; ele diz se a
      aplicação está resolvendo.
    \item Resistir à tentação de modelar todos os domínios antes de entregar o primeiro.
  \end{itemize}
}

\pendente{escrever}

% ---------------------------------------------------------------------------
\section{Uma arquitetura mínima}
\label{sec:arquitetura-minima}

\rascunho{
  \begin{itemize}
    \item As peças da Figura~\ref{fig:aplicacao-propria}: ingestão das fontes (de
      preferência a partir da camada prata, capítulo~\ref{cap:medalhao}), motor de
      conciliação, fila de revisão, registro mestre com a tabela de correspondência, e
      notificações.
    \item Três caminhos para cada par encontrado: alta confiança vai direto para o
      registro mestre; dúvida vai para revisão; baixa confiança é descartada.
    \item As decisões do steward voltam para o motor como regras e exemplos.
    \item Tecnologia: o mínimo que o time já domina. Um banco de dados, Python, um
      agendador e um canal de mensagens resolvem o começo.
  \end{itemize}
}

\pendente{escrever}

\begin{figure}[htbp]
  \centering
  % Arquitetura minima de uma aplicacao propria de master data e governanca (cap. 15).
\begin{tikzpicture}[diagrama, x=1mm, y=1mm]
  % fontes e ingestao
  \foreach \y/\nome/\id in {8/ERP/f1, 0/CRM/f2, -8/Planilhas/f3}
    \node[caixa, minimum width=14mm, minimum height=6mm] (\id) at (7,\y) {\nome};
  \node[moldura, fit=(f1)(f3), label={[rotulo]above:Fontes}] {};
  \node[caixa, minimum width=17mm, minimum height=10mm] (ing) at (29,0) {Ingestão\\{\scriptsize (prata)}};
  \foreach \id in {f1,f2,f3} \draw[seta] (\id.east) -- (ing.west);

  % motor de conciliacao
  \node[destaquecaixa, minimum width=24mm, minimum height=18mm] (mot) at (54,0)
    {Motor de\\conciliação\\[0.5mm]{\scriptsize regras,}\\{\scriptsize similaridade, IA}};
  \draw[seta] (ing) -- (mot);

  % saidas
  \node[bancodestaque, minimum width=24mm, minimum height=13mm] (reg) at (102,9)
    {\scriptsize Registro mestre e\\\scriptsize correspondência};
  \node[caixa, minimum width=20mm, minimum height=9mm] (fila) at (81,-10) {Fila de\\revisão};
  \node[caixa, minimum width=18mm, minimum height=9mm] (stew) at (104,-10) {Steward\\revisa};
  \node[caixa, draw=cinza, dashed, text=cinza, minimum width=20mm, minimum height=8mm] (notif) at (81,-24)
    {\scriptsize Mensagem\\\scriptsize aos donos};
  \draw[seta] (mot.east) ++(0,3) -- node[rotulo, above, sloped] {alta confiança} (reg.west);
  \draw[seta] (mot.east) ++(0,-3) -- node[rotulo, pos=0.55, below left=0.5mm] {dúvida} (fila.west);
  \draw[seta] (fila) -- (stew);
  \draw[seta] (stew) -- (reg);
  \draw[seta, draw=cinza] (fila) -- (notif);
  \node[rotulo, anchor=south] at (reg.north) {lido por BI, sistemas e IA};

  % aprendizado
  \draw[seta, dashed, draw=cinza] (stew.south) -- ++(0,-19) -| (mot.south);
  \node[rotulo, anchor=north] at (68,-33) {as decisões do steward viram regras e exemplos para o motor};
\end{tikzpicture}

  \caption{Uma arquitetura mínima para uma aplicação própria de master data e
    governança.}
  \label{fig:aplicacao-propria}
  \fonte{elaborado pelo autor.}
\end{figure}

% ---------------------------------------------------------------------------
\section{O motor de conciliação}
\label{sec:motor-conciliacao}

\rascunho{
  \begin{itemize}
    \item Em camadas, da mais barata para a mais cara: primeiro identificadores fortes
      (CNPJ, código do produto), depois regras de normalização e similaridade de texto,
      e só para o que sobrar, uma sugestão da IA com explicação.
    \item A base teórica é a conciliação probabilística de registros
      \parencite{fellegi1969,christen2012}. Existem bibliotecas de código aberto que a
      implementam. \pendente{conferir e citar bibliotecas atuais, como Splink e RapidFuzz}
    \item O trecho abaixo mostra a ideia mais simples, só com a biblioteca padrão do
      Python: normalizar os nomes e comparar a similaridade. Num sistema real, cada par
      vira um item na fila de revisão e uma mensagem para o dono do domínio.
    \item Os limiares (aqui, 0,85) se calibram com uma amostra revisada pelo negócio,
      e se reveem com as decisões do steward.
  \end{itemize}
}

\pendente{escrever}

\begin{python}
import unicodedata
from difflib import SequenceMatcher
from itertools import combinations

SUFIXOS = {"ltda", "sa", "me", "eireli"}

def normalizar(nome: str) -> str:
    """Minúsculas, sem acento, pontuação ou sufixo."""
    nome = unicodedata.normalize("NFKD", nome)
    nome = nome.encode("ascii", "ignore").decode().lower()
    nome = "".join(c for c in nome if c.isalnum() or c == " ")
    partes = nome.split()
    return " ".join(p for p in partes if p not in SUFIXOS)

def pares_suspeitos(nomes, limiar=0.85):
    """Pares de nomes que parecem a mesma entidade."""
    for a, b in combinations(sorted(set(nomes)), 2):
        na, nb = normalizar(a), normalizar(b)
        nota = SequenceMatcher(None, na, nb).ratio()
        if nota >= limiar:
            yield a, b, round(nota, 2)

clientes = ["AGRO SILVA LTDA", "Agro Silva", "Agrosilva",
            "Cerealista Norte S.A."]
for a, b, nota in pares_suspeitos(clientes):
    print(f"Revisar: {a!r} x {b!r} ({nota})")
\end{python}

\begin{terminal}
Revisar: 'AGRO SILVA LTDA' x 'Agro Silva' (1.0)
Revisar: 'AGRO SILVA LTDA' x 'Agrosilva' (0.95)
Revisar: 'Agro Silva' x 'Agrosilva' (0.95)
\end{terminal}

% ---------------------------------------------------------------------------
\section{Revisão humana e notificações}
\label{sec:revisao-notificacao}

\rascunho{
  \begin{itemize}
    \item A fila de revisão é o coração da aplicação: cada item mostra os registros, a
      nota, o motivo da sugestão e os botões aprovar, corrigir, rejeitar. Uma interface
      simples basta. \pendente{decidir se cito alguma ferramenta para interfaces internas}
    \item Notificação para quem precisa saber, pelo canal que a pessoa já usa (e-mail,
      chat corporativo), com o link para revisar. Mensagem demais vira ruído: agrupar por
      dia ou por domínio.
    \item Prazo combinado para a revisão, e um indicador de itens parados.
  \end{itemize}
}

\pendente{escrever}

% ---------------------------------------------------------------------------
\section{Histórico, permissões e privacidade}
\label{sec:auditoria-privacidade}

\rascunho{
  \begin{itemize}
    \item Toda mudança no registro mestre fica registrada: o que mudou, quem decidiu,
      quando, e a partir de qual sugestão. Sem isso não há governança, só um banco editável.
    \item Permissões por domínio: quem vê, quem sugere, quem aprova.
    \item Cadastro de cliente tem dado pessoal: mascarar antes de mandar qualquer coisa
      para um modelo de IA externo (ver \veja{sec:seguranca}).
  \end{itemize}
}

\pendente{escrever}

% ---------------------------------------------------------------------------
\section{Operar como produto}
\label{sec:operar-produto}

\rascunho{
  \begin{itemize}
    \item A aplicação é um produto de dados como os outros: código versionado, testes,
      entrega contínua e monitoramento (Parte~\ref{parte:devops}).
    \item Métricas próprias: porcentagem conciliada automaticamente, tempo médio de
      revisão, itens parados, duplicados encontrados por carga, custo por execução
      (inclusive em tokens, capítulo~\ref{cap:custo-consumo}).
    \item Um dono da aplicação, separado dos donos dos domínios.
  \end{itemize}
}

\pendente{escrever}

% ---------------------------------------------------------------------------
\section{Lista de verificação}
\label{sec:checklist-aplicacao}

\rascunho{
  A Tabela~\ref{tab:checklist-aplicacao} fecha o capítulo como referência rápida para
  quem vai começar.
}

\begin{table}[htbp]
  \centering
  \caption{Perguntas para responder antes de construir a própria aplicação.}
  \label{tab:checklist-aplicacao}
  \small
  \begin{tabularx}{\textwidth}{@{}>{\raggedright\arraybackslash}p{0.22\textwidth}
      >{\raggedright\arraybackslash}X@{}}
    \toprule
    Tema & Pergunta \\
    \midrule
    Domínio e dono & Qual domínio vem primeiro, e quem decide as regras dele? \\
    Fontes & De onde vêm os registros, e com que frequência mudam? \\
    Identidade & O que faz dois registros serem o mesmo? Existe um identificador forte? \\
    Limiares & Acima de que nota aprova sozinho? Abaixo de qual descarta? \\
    Revisão & Quem revisa, em quanto tempo e por qual canal? \\
    Notificação & Quem precisa saber de cada mudança, e com que frequência? \\
    Histórico & Toda mudança fica registrada, com quem decidiu e quando? \\
    Privacidade & O que pode ir para um modelo de IA, e o que precisa ser mascarado? \\
    Testes & Unicidade, correspondência e regras de negócio rodam a cada carga? \\
    Custo & Quanto custa cada execução, inclusive em tokens? \\
    Manutenção & Quem cuida da aplicação daqui a dois anos? \\
    Saída & Se um dia migrar para uma ferramenta de mercado, os dados saem fácil? \\
    \bottomrule
  \end{tabularx}
  \fonte{elaborado pelo autor.}
\end{table}

% ---------------------------------------------------------------------------
\begin{leituras}
  \leitura{christen2012}{Para aprofundar o motor de conciliação: normalização,
    bloqueio, comparação e avaliação dos pares encontrados.}
  \leitura{loshin2009}{Para os componentes de uma solução completa de master data, útil
    para saber o que se está deixando de fora ao começar pequeno.}
\end{leituras}
